% Standalone article. Build with: latexmk -pdf article.tex
\documentclass[11pt]{article}
\usepackage[T1]{fontenc}
\usepackage[utf8]{inputenc}
\usepackage{lmodern}
\usepackage[a4paper,margin=25mm,headheight=15pt]{geometry}
\usepackage{amsmath,amssymb,amsthm,mathtools}
\usepackage{microtype,booktabs,longtable,array,enumitem}
\usepackage[dvipsnames]{xcolor}
\usepackage{fancyhdr}
\usepackage{xurl}
\usepackage[colorlinks=true,linkcolor=MidnightBlue,citecolor=MidnightBlue,urlcolor=MidnightBlue]{hyperref}
\usepackage[nameinlink,noabbrev]{cleveref}
\hypersetup{pdftitle={Intrinsic Differential Calculus on the Surreal and Surcomplex Numbers},pdfauthor={Research exposition prepared for the Surreal documentation},pdfsubject={Berarducci--Mantova derivation, coherent calculus, bounded phases, and differential equations}}
\pagestyle{fancy}
\fancyhf{}
\fancyhead[L]{\small Intrinsic surreal and surcomplex calculus}
\fancyhead[R]{\small\thepage}
\renewcommand{\headrulewidth}{0.3pt}
\setlength{\parskip}{3pt}
\setlength{\parindent}{1.2em}
\setcounter{tocdepth}{2}
\emergencystretch=2em
\numberwithin{equation}{section}
\newtheorem{theorem}{Theorem}[section]
\newtheorem{proposition}[theorem]{Proposition}
\newtheorem{lemma}[theorem]{Lemma}
\newtheorem{corollary}[theorem]{Corollary}
\theoremstyle{definition}
\newtheorem{definition}[theorem]{Definition}
\newtheorem{example}[theorem]{Example}
\newtheorem{remark}[theorem]{Remark}
\newcommand{\No}{\mathbf{No}}
\newcommand{\R}{\mathbb R}
\newcommand{\C}{\mathbb C}
\newcommand{\Q}{\mathbb Q}
\newcommand{\Z}{\mathbb Z}
\newcommand{\N}{\mathbb N}
\newcommand{\K}{K}
\newcommand{\Lfield}{L}
\newcommand{\OO}{\mathcal O}
\newcommand{\mm}{\mathfrak m}
\newcommand{\II}{\mathcal I}
\newcommand{\HH}{\mathcal H}
\newcommand{\DD}{\partial}
\newcommand{\Dc}{\partial_{\mathrm{coef}}}
\newcommand{\Dz}{\partial_Z}
\newcommand{\Prim}{\mathcal A}
\newcommand{\Exp}{\operatorname{Exp}}
\newcommand{\Log}{\operatorname{Log}}
\newcommand{\ct}{\operatorname{ct}}
\newcommand{\st}{\operatorname{st}}
\newcommand{\supp}{\operatorname{supp}}
\newcommand{\ev}{\operatorname{ev}}
\newcommand{\cis}{\operatorname{cis}}
\newcommand{\im}{\operatorname{im}}
\newcommand{\re}{\operatorname{Re}}
\newcommand{\imag}{\operatorname{Im}}
\newcommand{\id}{\operatorname{id}}
\newcommand{\val}{v}
\newcommand{\repo}{\texttt{VladimirReshetnikov/Surreal}}
\newcommand{\snapshot}{\texttt{e260237db9b71da8b74a0c13c8e6355119091100}}
\newcommand{\pathtext}[1]{\nolinkurl{#1}}
\newcommand{\assertion}[1]{\par\smallskip\noindent\textbf{#1}\quad}
\begin{document}
\begin{center}
{\LARGE\bfseries Intrinsic Differential Calculus\\[4pt]
on the Surreal and Surcomplex Numbers}\par
\vspace{10pt}
{\large Strong derivations, coherent analytic families,\\
bounded phases, and linear differential equations}\par
\vspace{12pt}
{Research exposition prepared for the \repo\ documentation}\par
{21 September 2026}
\end{center}
\vspace{5pt}
\begin{abstract}
The repository develops an extensive function-variable calculus on
$\No[i]$, but explicitly leaves the choice of an intrinsic derivation outside
its trigonometric theory. This article supplies the missing bridge. We fix the
Berarducci--Mantova derivation $\DD$ on $K=\No$, distinguish it from
$\partial_Z$, and prove a total-derivative formula for common-domain
Hahn-coherent analytic families. Differentiation preserves that coherence,
although it may require a larger monomial workspace.

The central application is an exact obstruction to complex exponential
integration. Let $\OO$ be the finite surreals, $\mm$ the infinitesimals, and
$\II=\DD(\OO)=\DD(\mm)$. We prove that $\II$ is a nonzero convex ideal of
$\OO$, strictly smaller than $\mm$, and that
\[
 \left\{\frac{\DD y}{y}:y\in K[i]^\times\right\}=K+i\II.
\]
Consequently $\DD y=ay$ has a nonzero surcomplex solution precisely when the
imaginary part of $a$ has a finite surreal primitive. The canonical
exponential on $K+i\OO$ is compatible with $\DD$, and this is its maximal
possible domain even for pointwise differential compatibility. In particular,
$\DD y=iy$ and $\DD^2y+y=0$ have no nonzero solutions in $K[i]$.

We also classify the kernels of all complex constant-coefficient differential
operators and constant matrix systems, construct exact Laurent-series
resolvents, and provide a set-sized differential-workspace theorem. A
reproducible symbolic verifier checks finite identities and coefficient
recurrences; it is not a formal verification of the general proofs. The deep
existence and model-theoretic results are explicitly attributed to the
literature. No priority claim is made for the consequences developed here.
\end{abstract}
\vfill
{\small\noindent Repository snapshot: \snapshot.\\
Scope: a documentation-gap article, not a proof audit of the whole repository.}
\newpage
\begingroup
\small
\setlength{\parskip}{0pt}
\setlength{\baselineskip}{11pt}
\tableofcontents
\endgroup
\newpage

\section{The documentation gap and the mathematical route}
\label{sec:gap}

\subsection{What was inspected}
The audit was made against the pinned commit \snapshot, rather than against a
moving branch. The repository catalogue identifies fifteen report packages:
six about surreal arithmetic and constructions, seven about surcomplex
analysis and geometry, and two about foundations and computation
\cite{RepoMap,RepoManifest}. The principal evidence for the present choice
is the scope boundary in the trigonometry report: that report develops
finite-angle phases and global phase extensions, but expressly does not
choose a derivation on $\No$ or classify solutions of a differential equation
\cite{RepoTrig}. Its omission is deliberate, not a defect in its stated
results.

The analysis report distinguishes germ-analytic functions, common-domain
coherent sections, and ordinary holomorphic lifts. Its derivatives and
primitive theorems belong to these function classes \cite{RepoAnalysis}.
The computer-algebra report mentions transseries and the
Berarducci--Mantova literature in its representation and capability discussion
\cite{RepoCAS}. A bibliographic mention, however, does not supply an interface
between intrinsic differentiation of coefficients and differentiation in an
analytic variable.

\begin{center}
\begin{tabular}{@{}>{\raggedright\arraybackslash}p{.20\textwidth}>{\raggedright\arraybackslash}p{.34\textwidth}>{\raggedright\arraybackslash}p{.38\textwidth}@{}}
\toprule
Existing layer & Already supplied & The bridge developed here \\
\midrule
Analysis & Common-domain Hahn families and $Z$-derivatives & A coefficient derivation and the exact total-derivative rule \\
Trigonometry & Finite angles, unit-circle phases, phase extensions & Differential compatibility and a bounded-primitive obstruction \\
Computer algebra & Exact representation tiers and support contracts & Derivation-aware parents, Laurent resolvents, and verifiable residuals \\
Foundations & Class/set distinctions and local workspaces & Closure under differentiation and chosen integration, not just field operations \\
\bottomrule
\end{tabular}
\end{center}

This is a targeted scope audit, based on the catalogue and the relevant
package documentation. It is not a claim that the word ``derivation'' is
absent from every source manuscript, nor that every theorem in the existing
collection has been checked. The concrete gap is a connected development of
\emph{intrinsic differential algebra and its interaction with the already
present analytic and phase constructions}.

\subsection{Three results to keep in view}
First, for a coherent family $F(Z)$ evaluated at a finite surcomplex point
$z$, there are two contributions to the derivative:
\begin{equation}\label{eq:introchain}
 \DD\bigl(\ev_zF\bigr)
 =\ev_z(\Dc F)+\ev_z(\Dz F)\,\DD z.
\end{equation}
The first term is indispensable: the coefficients of $F$ are surreal
numbers and need not be constants of $\DD$.

Second, a nonzero surcomplex number has a canonical infinitesimal phase
relative to its ordinary unit-circle direction. That phase has an intrinsic
derivative, but not every surreal number is the derivative of a finite
phase. This gives the exact logarithmic-derivative image
$K+i\DD(\mm)$.

Third, absence of oscillatory homogeneous solutions is compatible with
excellent inhomogeneous solvability. In $\C((\omega^{-1}))$, for example,
$\DD-i$ is bijective, even though its homogeneous equation has no nonzero
solution anywhere in $\No[i]$. These are different mathematical properties.

\subsection{What is assumed, and what is proved}
We take the construction of the surreal ordered exponential field, its Hahn
normal forms, and the Berarducci--Mantova existence theorem as external
inputs. The elementary-model statement for transseries is also cited, not
reproved. The other main results of this article are proved from the stated
inputs. Their purpose is to complete the documentation interface, not to
announce a newly solved published problem. In particular, ordinary
constant-coefficient differential algebra and formal resolvents have a long
history; their specialization here is exposition, not a priority claim.

\section{Algebraic setting and the intrinsic derivation}
\label{sec:setting}

\subsection{Classes, constants, and finite parts}
Work in a classical class theory such as NBG, with the usual class
constructions of the surreals. Put
\[
 K=\No,\qquad L=K[i]=\{a+ib:a,b\in K\},\qquad i^2=-1.
\]
$K$ is real closed, and $L$ is algebraically closed. Statements below about
linear algebra involve only finite matrices or explicitly specified
set-sized subfields. No summation over a proper class is used.

For $a,b\in K$, write $a\prec b$ when $b\ne0$ and
$|a|<|b|/n$ for every positive integer $n$. Define
\[
 \OO=\{a\in K:|a|<n\text{ for some }n\in\N\},\qquad
 \mm=\{a\in K:|a|<1/n\text{ for every }n\ge1\}.
\]
Then $\OO=\R\oplus\mm$ additively. The projection
$\st:\OO\to\R$ is a ring homomorphism with kernel $\mm$.
On $L$ use
\[
 \overline{a+ib}=a-ib,\qquad |a+ib|=\sqrt{a^2+b^2}.
\]
The finite and infinitesimal subclasses are
$\OO_L=\OO+i\OO$ and $\mm_L=\mm+i\mm$; standard part is taken
coordinatewise.

A Conway normal form is a set-supported sum
\[
 x=\sum_{\alpha\in S}r_\alpha\omega^\alpha,
\]
with real coefficients and reverse-well-ordered support $S$. Here and below,
$\omega^\alpha$ denotes the Conway monomial map. The coefficient of the
monomial $1=\omega^0$ defines an additive real-linear map
\begin{equation}\label{eq:ct}
 \ct:K\longrightarrow\R.
\end{equation}
It agrees with $\st$ on $\OO$, but is not multiplicative on $K$:
$\ct(\omega)\ct(\omega^{-1})=0$, whereas
$\ct(\omega\omega^{-1})=1$.

\subsection{The external existence theorem}
\begin{theorem}[Berarducci--Mantova; external input]\label{thm:BM}
There is a derivation $\DD=\partial_{\mathrm{BM}}:K\to K$ with
\begin{align}
 \DD(xy)&=x\DD y+y\DD x, & \ker\DD&=\R,\label{eq:BM1}\\
 \DD\!\left(\sum_jx_j\right)&=\sum_j\DD x_j,
 &\DD(\exp x)&=(\exp x)\DD x,\label{eq:BM2}\\
 x>\R&\Longrightarrow\DD x>0,
 &\DD\omega&=1.\label{eq:BM3}
\end{align}
In \eqref{eq:BM2}, every summable set-indexed family has a summable family of
derivatives. Moreover, $\DD$ is surjective.
\end{theorem}

These are the properties of the distinguished derivation constructed in
\cite[Theorems A and B]{BM}; the published version is
\cite{BMpub}. They imply that $K$ is a Liouville closed $H$-field.
Existence of several surreal derivations is part of the same work; the
listed structural ideas are not an assertion that the underlying field has
only one possible derivation.

\assertion{Construction perspective.}
One cannot define $\DD$ simply by recursively differentiating an exponential
expression for every number. Log-atomic numbers require extra initial data.
The construction specifies a pre-derivation there and extends it through
summable path derivatives. The summability proof and surjectivity proof are
substantial inputs, not consequences of the product rule alone
\cite[\S\S6--7]{BM}. This article uses their conclusions rather than a
pretended implementation of that construction.

\subsection{First consequences}
All real numbers are constants, including real parameters that may occur in
an exponent. Thus
\begin{equation}\label{eq:basicderivatives}
 \DD(\omega^q)=q\omega^{q-1}\quad(q\in\R),\qquad
 \DD(\log\omega)=\omega^{-1},\qquad
 \DD(\exp\omega)=\exp\omega.
\end{equation}
For rational $q$ the first formula follows from algebraic differentiation;
for real $q$ use $\omega^q=\exp(q\log\omega)$ for real exponents and
exponential compatibility. The logarithm identity follows by
differentiating $\exp(\log x)=x$ for $x>0$.

For $\ell_0=\omega$ and $\ell_{n+1}=\log\ell_n$, induction gives
\begin{equation}\label{eq:logtower}
 \DD\ell_n=\frac{1}{\ell_0\ell_1\cdots\ell_{n-1}}
 \quad(n\ge1).
\end{equation}
All these $\ell_n$ are positive infinite. The empty product for $n=0$ is
understood as $1$.

\begin{remark}[Two different exponent constructions]
For an arbitrary surreal exponent $a$, the Conway monomial $\omega^a$ must
not be silently replaced by $\exp(a\log\omega)$. The latter expression has
derivative
\[
 \exp(a\log\omega)\left((\DD a)\log\omega+a/\omega\right),
\]
but this does not give a general rule for the Conway monomial map. A
symbolic implementation needs different constructors for the two operations.
\end{remark}

\subsection{Why this is not a derivative of a constant function}
As an element of $K$, $\omega$ satisfies $\DD\omega=1$. As the constant
function $Z\mapsto\omega$, it satisfies $\partial_Z\omega=0$. Similarly,
for the ordinary function $f(Z)=e^Z$ at a real point $r$,
\[
 f'(r)=e^r,\qquad \DD(f(r))=\DD(e^r)=0.
\]
The first derivative varies an argument; the second acts on an element of a
differential field. Neither equation contradicts the other.

The transseries embedding gives a useful interpretation: the natural
embedding sending the transseries variable to $\omega$ is elementary as an
embedding of ordered differential fields \cite[Theorem 1]{ADHsurreal}.
This explains the asymptotic intuition behind $\DD\omega=1$. It does not
identify every arbitrary surcomplex function with a transseries.

\section{Small derivatives and normalized integration}
\label{sec:integration}

\subsection{An order reversal on infinitesimals}
\begin{lemma}\label{lem:smallnegative}
If $0<\varepsilon\in\mm$, then $\DD\varepsilon<0$. Moreover
$\DD(\OO)\subseteq\mm$.
\end{lemma}
\begin{proof}
$1/\varepsilon$ is positive infinite, so the $H$-field axiom gives
\[
 0<\DD(1/\varepsilon)=-\frac{\DD\varepsilon}{\varepsilon^2}.
\]
Hence $\DD\varepsilon<0$. For each positive integer $n$,
$\omega+n\varepsilon$ is positive infinite. Its derivative is
$1+n\DD\varepsilon>0$. Therefore
$-1/n<\DD\varepsilon<0$ for every $n$, so
$\DD\varepsilon\in\mm$. Negative infinitesimals follow by additivity.
Every finite element is a real constant plus an infinitesimal.
\end{proof}

In particular, $\DD$ is strictly order-reversing on $\mm$: if
$\varepsilon_1<\varepsilon_2$, apply the lemma to their positive
infinitesimal difference. This order reversal is not a global property of
$\DD$; positive infinite elements have positive derivative.

\subsection{A fixed integration normalization}
\begin{proposition}\label{prop:A}
For each $b\in K$ there is a unique $\Prim(b)\in K$ such that
\[
 \DD\Prim(b)=b,\qquad \ct(\Prim(b))=0.
\]
The map $\Prim$ is real-linear and satisfies
\begin{equation}\label{eq:Aidentities}
 \DD\Prim=\id_K,\qquad \Prim\DD=\id_K-\ct.
\end{equation}
\end{proposition}
\begin{proof}
Choose any primitive $B$ of $b$, using surjectivity, and replace it by
$B-\ct(B)$. If two normalized primitives existed, their difference would
lie in $\ker\DD=\R$ and have constant coefficient zero, hence be zero.
Additivity and real homogeneity follow from uniqueness. Applying this to
$\DD x$ gives the second identity.
\end{proof}

The existence proof does not require a simultaneous arbitrary choice of
primitives: the normalization defines a unique value for each input.
We do not claim that $\Prim$ preserves arbitrary Hahn sums, is multiplicative,
or is an effective procedure for arbitrary surreal inputs. Its usefulness
here is semantic precision. For example,
\[
 \Prim(1)=\omega,\qquad
 \Prim(\omega^{-1})=\log\omega-\ct(\log\omega).
\]
The second primitive is infinite regardless of its real constant term.

\subsection{The bounded-primitive ideal}
\begin{definition}\label{def:I}
The \emph{bounded-primitive ideal} is the additive subclass
\[
 \II:=\DD(\OO)=\DD(\mm)\subset K.
\]
The second equality follows because $\DD$ kills $\R$.
\end{definition}

\begin{theorem}\label{thm:I}
The class $\II$ has the following properties.
\begin{enumerate}[label=(\roman*),itemsep=2pt]
\item $\DD:\mm\to\II$ is a real-linear, order-reversing bijection.
\item $b\in\II$ if and only if $\Prim(b)\in\mm$, equivalently if and only if
some, and hence every, primitive of $b$ is finite.
\item $\II$ is convex in $K$ and is an ideal of the valuation ring $\OO$.
\item $0\ne\II\subsetneq\mm\subsetneq\OO$.
\item There is an exact cut description
\begin{equation}\label{eq:Icut}
 \II=\{b\in K: |b|<\DD f\text{ for every }f>\R\}.
\end{equation}
\end{enumerate}
\end{theorem}
\begin{proof}
The first two claims follow from \cref{lem:smallnegative,prop:A} and the
constant field. To prove convexity, suppose $0\le b\le d$ with
$d\in\II$, $d\ge0$. Write $d=-\DD\varepsilon$ with
$\varepsilon\in\mm$, $\varepsilon\ge0$, using order reversal. Let
$B'=b$. If $B$ were negative infinite, then $-B$ would be positive infinite
and $b<0$, a contradiction. If $B$ were positive infinite, then
$B+\varepsilon$ would be positive infinite but
\[
 \DD(B+\varepsilon)=b-d\le0,
\]
another contradiction. Thus $B$ is finite and $b\in\II$. Sign symmetry
proves convexity in general.

For $u\in\OO$ and $b\in\II$, choose a positive integer $n>|u|$.
Then $|ub|\le n|b|$, and convexity and real-linearity imply
$ub\in\II$. Thus $\II$ is an ideal of $\OO$.

The lemma already gives $\II\subseteq\mm$. It is nonzero because
$\DD(-\omega^{-1})=\omega^{-2}$. But $\omega^{-1}\notin\II$:
every primitive differs from $\log\omega$ by a real constant and is
infinite. This proves strictness.

For \eqref{eq:Icut}, if $b=\DD\varepsilon$ with
$\varepsilon\in\mm$, then $f+\varepsilon$ and $f-\varepsilon$ are
positive infinite for every $f>\R$. Their positive derivatives give
$\DD f>|b|$. Conversely, suppose $b\notin\II$, and choose $B'=b$.
Then $B$ is infinite. For $f=|B|/2$ the $H$-field axiom gives
$\DD|B|=|b|>0$, whence $\DD f=|b|/2$. The inequality required in
\eqref{eq:Icut} fails.
\end{proof}

The inverse $J:\II\to\mm$ of $\DD|_\mm$ is the unique infinitesimal
primitive. It equals $\Prim|_\II$. In contrast to a general primitive, it
has no residual additive constant ambiguity.

\begin{corollary}\label{cor:quotients}
Differentiation induces an order-preserving real-linear isomorphism of
ordered additive quotients
\[
 K/\OO\ \longrightarrow\ K/\II,\qquad x+\OO\longmapsto\DD x+\II.
\]
\end{corollary}
\begin{proof}
Surjectivity is inherited from $\DD$, and
$\DD^{-1}(\II)=\OO$. If $x+\OO$ is positive, then $x$ is positive
infinite. For any $d=\DD\varepsilon\in\II$, the element
$x-\varepsilon$ is positive infinite, so $\DD x>d$. Hence the induced
map preserves positivity. The converse follows by bijectivity and sign
symmetry.
\end{proof}

\subsection{Examples at increasingly slow scales}
For real $p$, integration gives
\[
 \DD\!\left(\frac{\omega^{1-p}}{1-p}\right)=\omega^{-p}
 \quad(p\ne1).
\]
Therefore
\begin{equation}\label{eq:powerthreshold}
 \omega^{-p}\in\II\quad\Longleftrightarrow\quad p>1.
\end{equation}
The endpoint $p=1$ fails because its primitive is $\log\omega$.

More generally, put
\[
 b_{n,p}=\frac{1}{\ell_0\ell_1\cdots\ell_{n-1}\ell_n^p}
 \quad(n\ge0),
\]
where the prefix product is empty for $n=0$. For $p\ne1$ a primitive is
$\ell_n^{1-p}/(1-p)$; for $p=1$ it is $\ell_{n+1}$. Thus
\begin{equation}\label{eq:logthreshold}
 b_{n,p}\in\II\quad\Longleftrightarrow\quad p>1.
\end{equation}
Being infinitesimal is not enough. In particular,
$1/(\omega\log\omega)$ is infinitesimal but not in $\II$.
These tests are exact algebraic statements about primitives, not claims
about real improper integrals.

\section{Complexification and infinitesimal formal calculus}
\label{sec:complexification}

\subsection{The extension of the derivation}
\begin{proposition}\label{prop:complexextension}
There is exactly one derivation on $L=K[i]$ extending $\DD$ on $K$. It is
\[
 \DD(a+ib)=\DD a+i\DD b.
\]
It commutes with conjugation, has constant field $\C$, is surjective, and
is strongly additive on set-indexed Hahn-summable families in $L$.
\end{proposition}
\begin{proof}
Differentiating $i^2=-1$ gives $2i\DD i=0$, so $\DD i=0$ in any
extension. The displayed formula is therefore forced, and direct expansion
verifies its product rule. The assertions about constants and surjectivity
follow coordinatewise. For strong additivity, take real and imaginary
parts of a summable family. Their derivative families are summable by
\cref{thm:BM}; a finite union of reverse-well-ordered supports is again
reverse well ordered, and each monomial still has finite multiplicity.
\end{proof}

The normalized primitive extends to
$\Prim_L(a+ib)=\Prim(a)+i\Prim(b)$. It is $\C$-linear, satisfies
$\DD\Prim_L=\id$, and has zero complex constant term. This is an
algebraically specified class operation, not a promise of an effective
integration algorithm for arbitrary input encodings.

For $z\ne0$, write
\[
 z^\dagger=\frac{\DD z}{z}.
\]
The logarithmic derivative satisfies
\begin{equation}\label{eq:logderivative}
 (zw)^\dagger=z^\dagger+w^\dagger,\qquad
 (z^{-1})^\dagger=-z^\dagger,\qquad
 \ker({}^\dagger)=\C^\times.
\end{equation}
This notation does not presume that a global surcomplex logarithm exists.

\begin{lemma}\label{lem:radial}
For $z\in L^\times$ and $r=|z|$, one has
\[
 \frac{\DD r}{r}=\re(z^\dagger).
\]
\end{lemma}
\begin{proof}
Differentiate $r^2=z\bar z$ and divide by $2r^2$:
$\DD r/r=(z^\dagger+\overline{z^\dagger})/2$.
\end{proof}

\subsection{Summability is a support condition}
A family $(x_j)_{j\in J}$ of Hahn series is \emph{summable} when $J$ is a
set, the union of their monomial supports is reverse well ordered, and each
monomial occurs in only finitely many of the $x_j$. The sum is then
coefficientwise. This is not convergence of partial sums in the fine
surreal topology.

We use the standard Hahn--Neumann support lemma: if $S$ is a
reverse-well-ordered set of monomials strictly smaller than $1$, the
multiplicative semigroup it generates is reverse well ordered; for any
fixed monomial, only finitely many products from $S$ occur, counting
multiplicities up to permutation. These are the summability facts used in
the formal-series calculus of \cite[\S3.2]{BMgerms}. They hold for set supports,
without a restriction to an Archimedean value group.

Consequently, for $h\in\mm_L$ and an arbitrary formal series
$f(X)=\sum_{n\ge0}c_nX^n\in\C[[X]]$, the family $(c_nh^n)_n$ is
summable. Its value is denoted by $f(h)$. No lower bound on an ordinary
radius of convergence is needed: factorially growing complex coefficients
remain coefficients at the same Hahn monomials.

\begin{proposition}[Infinitesimal chain rule]\label{prop:formalchain}
For $f\in\C[[X]]$ and $h\in\mm_L$,
\begin{equation}\label{eq:formalchain}
 \DD(f(h))=f_{\mathrm{formal}}'(h)\,\DD h.
\end{equation}
\end{proposition}
\begin{proof}
For each $n$, the product rule gives
$\DD(c_nh^n)=nc_nh^{n-1}\DD h$. Strong additivity applies to the
summable family $(c_nh^n)_n$. The family
$(nc_nh^{n-1})_{n\ge1}$ is separately summable by the support lemma,
and multiplication by the fixed number $\DD h$ preserves summability.
Hence its sum can be factored as stated. The case $h=0$ is immediate.
\end{proof}

The same support argument justifies addition, multiplication and formal
composition whenever the inner series has zero constant term. In
particular, define
\begin{equation}\label{eq:smallExpLog}
 \Exp(h)=\sum_{n\ge0}\frac{h^n}{n!},\qquad
 \Log(1+h)=\sum_{n\ge1}\frac{(-1)^{n+1}h^n}{n}.
\end{equation}
These are inverse group isomorphisms
\[
 (\mm_L,+)\ \underset{\Log}{\stackrel{\Exp}{\rightleftarrows}}\
 (1+\mm_L,\cdot).
\]
Formal identities give $\Exp(h+k)=\Exp(h)\Exp(k)$, and
\begin{equation}\label{eq:smallExpLogderivatives}
 \DD\Exp(h)=\Exp(h)\DD h,\qquad
 \DD\Log(1+h)=\frac{\DD h}{1+h}.
\end{equation}
Both functions commute with conjugation. On real infinitesimal inputs,
$\Exp$ agrees with the surreal exponential, and $\Log$ with its inverse
on $1+\mm$.

\begin{remark}[An important limit on the evaluation rule]
The rule $f(h)$ above assumes \emph{constant} coefficients
$c_n\in\C$. A formal series with arbitrary coefficients in $L$ need not
be summable at any given nonzero infinitesimal. For example,
$\sum_{n\ge0}\omega^{n^2}X^n$ is not summable at $X=\omega^{-1}$.
Its monomials $\omega^{n^2-n}$ are not reverse well ordered. One must
not export the infinitesimal chain rule to arbitrary coefficient streams
without a joint support certificate.
\end{remark}

\section{A differential interface for coherent surcomplex analysis}
\label{sec:coherent}

\subsection{One coefficient domain, two derivations}
Let $U\subseteq\C$ be an ordinary open set. Consider the common-domain
Hahn algebra whose elements have the form
\begin{equation}\label{eq:coherentF}
 F(Z)=\sum_{m\in S}f_m(Z)m,
 \qquad f_m\in\mathcal O(U),
\end{equation}
where $S$ is a set of surreal monomials, reverse well ordered. Zero
coefficient functions may be discarded. All the coefficient functions live
on the \emph{same} domain $U$. This is a common-domain construction of the
kind distinguished in the repository, not the radius-free algebra of
unrelated convergent germs \cite{RepoAnalysis}.

Its ordinary analytic-variable derivation is
\[
 \Dz F=\sum_{m\in S}f_m'(Z)m.
\]
Its intrinsic coefficient derivation should instead satisfy
\[
 \Dc F=\sum_{m\in S}f_m(Z)\DD m.
\]
The second display is meaningful only after checking that its expanded
coefficients remain a valid common-domain family.

\begin{theorem}[Preservation of common-domain coherence]
\label{thm:coefficientderivation}
The operation $\Dc$ is a well-defined derivation of the common-domain
Hahn algebra in \eqref{eq:coherentF}. It preserves the ordinary domain
$U$ and commutes with $\Dz$:
\begin{equation}\label{eq:commutingderivations}
 \Dc\Dz F=\Dz\Dc F.
\end{equation}
It need not preserve a previously chosen smaller monomial group.
\end{theorem}
\begin{proof}
The family $(m)_{m\in S}$ is summable. By strong additivity, the family
$(\DD m)_{m\in S}$ is summable. Write
\[
 \DD m=\sum_n d_{n,m}n,\qquad d_{n,m}\in\R.
\]
The union $T$ of these supports is reverse well ordered, and, for each
$n\in T$, only finitely many $m$ have $d_{n,m}\ne0$. Consequently
\begin{equation}\label{eq:coefficientmatrix}
 \Dc F=\sum_{n\in T}g_n(Z)n,
 \qquad g_n(Z)=\sum_{m\in S}d_{n,m}f_m(Z)
\end{equation}
is legitimate, and every inner sum is finite. Each $g_n$ is holomorphic
on the same $U$.

Linearity is immediate. The product rule follows by expanding the Hahn
product of two families: each coefficient receives finitely many
contributions, and $\DD(m_1m_2)=(\DD m_1)m_2+m_1\DD m_2$.
The support lemma permits the regrouping. Finally, ordinary differentiation
of the finite sum in \eqref{eq:coefficientmatrix} replaces each $f_m$ by
$f_m'$, giving \eqref{eq:commutingderivations}. The support of $\DD m$
can contain monomials outside a preselected subgroup, which accounts for
the last assertion.
\end{proof}

This proof makes the interface particularly concrete: the derivation is
represented by a column-indexed family $(\DD m)_m$ whose matrix is
\emph{row finite on every specified summable input support}. It does not
assert that every individual derivative has finite support, or that the
matrix has a finite global representation.

\subsection{Evaluation and total differentiation}
For $z\in\OO_L$ with $a=\st z\in U$, put $h=z-a\in\mm_L$ and
define
\begin{equation}\label{eq:evaluation}
 \ev_z F=
 \sum_{m\in S}\sum_{k\ge0}
     m\,\frac{f_m^{(k)}(a)}{k!}h^k.
\end{equation}
The family is jointly summable. Indeed, the Taylor coefficient at each
$(m,k)$ is an ordinary complex constant, and the supports are controlled
by $S$ times the semigroup generated by the support of $h$. The
Hahn--Neumann lemma gives both reverse well ordering and finite
multiplicity. This remains true when the ordinary sizes of the Taylor
coefficients grow without a common bound.

\begin{theorem}[Total intrinsic derivative]\label{thm:totalchain}
For $F$ and $z$ as above,
\begin{equation}\label{eq:totalchain}
 \DD(\ev_zF)=\ev_z(\Dc F)+\ev_z(\Dz F)\,\DD z.
\end{equation}
\end{theorem}
\begin{proof}
Apply strong additivity to the double family in \eqref{eq:evaluation}.
Since $a\in\C$, all numbers $f_m^{(k)}(a)/k!$ are constants of
$\DD$, and $\DD h=\DD z$. The product rule splits the derivative
of each term into
\[
 (\DD m)\frac{f_m^{(k)}(a)}{k!}h^k
 \quad\text{and}\quad
 m\frac{f_m^{(k)}(a)}{(k-1)!}h^{k-1}\DD z
 \quad(k\ge1).
\]
Each of the two families is separately summable. For the first, use the
summability of $(\DD m)_m$ and the same infinitesimal semigroup. For the
second, use the support certificate for $\Dz F$ and multiply by the
fixed number $\DD z$. The first sum is the evaluation of
\eqref{eq:coefficientmatrix}; reindexing the second gives
$\ev_z(\Dz F)\DD z$.
\end{proof}

An ordinary holomorphic lift $f^\#$ has constant coefficients with respect
to $\DD$, so the formula reduces to
$\DD(f^\#(z))=(f')^\#(z)\DD z$. A general coherent family has the
additional coefficient term. This distinction is independent of whether
one can write its evaluated value as a short expression.

\begin{example}[A moving argument and moving coefficients]
Let $t=\omega^{-1}$, $F(Z)=\omega Z^2+tZ$, and $z=1+t$. Then
\[
 \Dc F=Z^2-t^2Z,\qquad \Dz F=2\omega Z+t,\qquad \DD z=-t^2.
\]
The evaluated value is
\[
 F(z)=\omega+2+2t+t^2,
 \qquad \DD(F(z))=1-2t^2-2t^3.
\]
Substituting in \eqref{eq:totalchain} gives exactly the same expression.
Keeping only the argument derivative would miss the leading constant $1$.
The simpler family $F(Z)=\omega$ already disproves that omission.
\end{example}

\subsection{Fixed contours and moving poles}
Suppose a fixed ordinary piecewise smooth contour $\gamma$ lies in $U$.
Define its coefficientwise integral by
\[
 \int_\gamma F(Z)\,dZ
   =\sum_{m\in S}\left(\int_\gamma f_m(Z)\,dZ\right)m.
\]
The finite sums in \eqref{eq:coefficientmatrix} show directly that
\begin{equation}\label{eq:contourcommutation}
 \DD\!\left(\int_\gamma F\,dZ\right)
       =\int_\gamma\Dc F\,dZ.
\end{equation}
The contour is fixed and ordinary, and the integral on either side is a
coefficientwise functional. No surcomplex-valued Riemann integral is
being introduced. The same argument applies to a coefficientwise residue
at a fixed ordinary centre whenever all coefficients are meromorphic on
one punctured neighbourhood: residue is a complex-linear functional, and
each new coefficient is a finite linear combination.

At an actual moving pole, the coefficient term cannot be ignored. For
$a\in L$ and the rational expression $F(Z)=1/(Z-a)$,
\begin{equation}\label{eq:movingpole}
 \Dc F=\frac{\DD a}{(Z-a)^2},\qquad
 \Dz F=-\frac{1}{(Z-a)^2},\qquad
 \DD\frac1{z-a}=-\frac{\DD z-\DD a}{(z-a)^2}.
\end{equation}
These rational identities hold whenever the denominators are nonzero;
they do not need a coherent expansion. The fixed-centre residue in
\eqref{eq:contourcommutation} must not be silently replaced by an actual
residue at $a$. The repository already distinguishes those residues;
intrinsic differentiation supplies another reason to retain the distinction.

\section{Canonical infinitesimal phase and exponential integration}
\label{sec:phase}

\subsection{A polar form with no branch choice}
The ordinary unit circle is
$\mathbb T(\C)=\{c\in\C:|c|=1\}$. Its surreal counterpart is
$\mathbb T(K)=\{u\in L:|u|=1\}$.

\begin{theorem}[Infinitesimal polar decomposition]\label{thm:polar}
Every $z\in L^\times$ has a unique expression
\begin{equation}\label{eq:polar}
 z=r\,c\,\Exp(i\theta),
 \qquad r\in K_{>0},\quad c\in\mathbb T(\C),\quad\theta\in\mm.
\end{equation}
Here $r=|z|$, $c=\st(z/|z|)$, and
$i\theta=\Log\bigl(z/(rc)\bigr)$. In particular,
\[
 \mathbb T(K)\cong\mathbb T(\C)\times(\mm,+)
\]
as abelian groups.
\end{theorem}
\begin{proof}
Set $r=|z|$ and $u=z/r$. Both coordinates of $u$ are finite, and
$u\bar u=1$. Taking standard parts gives $c\bar c=1$ for
$c=\st u$, so $c\ne0$. Then $w=u/c$ belongs to $1+\mm_L$ and
satisfies $w\bar w=1$. Put $q=\Log w$. Conjugation and the formal
logarithm identities give
\[
 q+\bar q=\Log(w\bar w)=0.
\]
Thus $q=i\theta$ for a unique real infinitesimal $\theta$.
Exponentiating gives \eqref{eq:polar}. Conversely, that expression has
modulus $r$ because $\Exp(i\theta)\Exp(-i\theta)=1$, has ordinary
unit direction $c$, and has logarithm $i\theta$ after division by $rc$.
These determine all three quantities uniquely. Multiplication multiplies
$r,c$ and adds $\theta$, proving the group statement.
\end{proof}

The unit-circle decomposition is closely related to the finite-angle
theorem already present in the repository \cite{RepoTrig}. What matters
here is that its last factor has \emph{a unique infinitesimal parameter},
so its intrinsic derivative has no branch ambiguity.

\begin{corollary}\label{cor:polarlogderivative}
For the polar decomposition \eqref{eq:polar},
\begin{equation}\label{eq:polarlogderivative}
 z^\dagger=\DD\log r+i\DD\theta.
\end{equation}
\end{corollary}
\begin{proof}
The ordinary number $c$ is a differential constant. Differentiate the
three factors using exponential compatibility and
\eqref{eq:smallExpLogderivatives}, then divide by $z$.
\end{proof}

\subsection{The exact logarithmic-derivative image}
\begin{theorem}[Surcomplex exponential-integrability criterion]
\label{thm:logderivativeimage}
The image of the logarithmic derivative on $L^\times$ is exactly
\begin{equation}\label{eq:logderivativeimage}
 (L^\times)^\dagger=K+i\II,
 \qquad \II=\DD(\OO)=\DD(\mm).
\end{equation}
For $a,b\in K$, the equation
\begin{equation}\label{eq:firstorder}
 \DD y=(a+ib)y
\end{equation}
has a nonzero solution in $L$ if and only if $b\in\II$. If this
condition holds, all its solutions are
\begin{equation}\label{eq:allfirstordersolutions}
 y=c\,\exp(\Prim(a))\,\Exp(iJ(b)),\qquad c\in\C.
\end{equation}
\end{theorem}
\begin{proof}
For necessity, \eqref{eq:polarlogderivative} shows that the imaginary
part of every logarithmic derivative is the derivative of an
infinitesimal, hence belongs to $\II$.

For sufficiency, $\Prim(a)$ is a primitive of $a$, and $J(b)\in\mm$
is the unique infinitesimal primitive of $b$. The product
$y_0=\exp(\Prim(a))\Exp(iJ(b))$ is nonzero and has logarithmic
derivative $a+ib$. Finally, for two nonzero solutions,
$\DD(y/y_0)=0$, so $y/y_0\in\C^\times$. Including $c=0$ gives
all solutions in \eqref{eq:allfirstordersolutions}.
\end{proof}

The real component is unrestricted because $K$ is Liouville closed. The
imaginary component is restricted not merely by smallness, but by the
much sharper finite-primitive condition. Algebraic closure of $L$ supplies
roots of polynomials in its elements; it does not supply roots of
arbitrary differential equations.

\begin{example}[The threshold is sharper than infinitesimality]
For real $p$, the equation
\[
 \DD y=i\omega^{-p}y
\]
has a nonzero solution if and only if $p>1$. For $p>1$ one may take
\[
 y=\Exp\!\left(\frac{i\omega^{1-p}}{1-p}\right).
\]
For $p=1$, a real primitive of the angular velocity is $\log\omega$,
which is infinite; for $p<1$, its primitive is also infinite. The
logarithmic-scale family \eqref{eq:logthreshold} gives arbitrarily slower
endpoint failures, for example $b=1/(\omega\log\omega)$.
\end{example}

There is a concise quotient formulation. The additive homomorphism
\[
 L\longrightarrow K/\II,\qquad a+ib\longmapsto b+\II
\]
is surjective and has kernel $(L^\times)^\dagger$. Therefore
\begin{equation}\label{eq:logderivativequotient}
 L/(L^\times)^\dagger\cong K/\II
\end{equation}
as additive real-vector-space quotients. These are not quotient fields.
The class $b+\II$ is the complete obstruction to exponential integration.

\section{The maximal differential exponential strip}
\label{sec:strip}

\subsection{Construction and kernel}
For a finite real $\theta=s+\varepsilon$, with $s=\st\theta\in\R$
and $\varepsilon\in\mm$, define
\[
 \cis(\theta)=e^{is}\Exp(i\varepsilon).
\]
This is an additive-to-multiplicative homomorphism on $\OO$. Its kernel
is $2\pi\Z$: standard part first forces $s\in2\pi\Z$, then
injectivity of the small exponential forces $\varepsilon=0$.

Let
\[
 \mathcal S=K+i\OO.
\]
It is an additive subgroup and a real-vector subspace of $L$, not a
subring: $\omega\in\mathcal S$ and $i\in\mathcal S$, but
$i\omega\notin\mathcal S$. Define
\begin{equation}\label{eq:stripExp}
 E(a+i\theta)=\exp(a)\cis(\theta),
 \qquad a\in K,\quad\theta\in\OO.
\end{equation}

\begin{theorem}\label{thm:stripexp}
The map $E:(\mathcal S,+)\to(L^\times,\cdot)$ is a surjective
homomorphism with kernel $2\pi i\Z$. It satisfies
\begin{equation}\label{eq:stripcompatibility}
 \DD E(w)=E(w)\DD w\qquad(w\in\mathcal S).
\end{equation}
Consequently
$\mathcal S/(2\pi i\Z)\cong L^\times$ as abelian groups.
\end{theorem}
\begin{proof}
Multiplicativity follows from the exponential and finite-phase group
laws. If $E(a+i\theta)=1$, taking moduli gives $\exp a=1$, hence
$a=0$, and the phase kernel gives $\theta\in2\pi\Z$. Conversely,
these elements are in the kernel.

For surjectivity, write $z=r c\Exp(i\varepsilon)$ by
\cref{thm:polar}. Choose an ordinary real $s$ with $e^{is}=c$; then
$z=E(\log r+i(s+\varepsilon))$. This uses the ordinary circle only,
not a branch on the entire surcomplex plane.

For compatibility, differentiate \eqref{eq:stripExp}. The ordinary
number $s=\st\theta$ is a constant, and
$\DD(\theta-s)=\DD\theta$. Therefore
$\DD\cis(\theta)=i\cis(\theta)\DD\theta$, which proves
\eqref{eq:stripcompatibility}.
\end{proof}

Every nonzero $z$ thus has strip logarithms, all differing by an ordinary
period $2\pi i n$. All such logarithms have the same intrinsic
derivative, namely $z^\dagger$. This derivative-level conclusion is
canonical even when a chosen logarithm branch is not.

\subsection{A pointwise maximality theorem}
\begin{theorem}[Maximal domain of differential compatibility]
\label{thm:maximalstrip}
For $w=a+iB\in L$, there exists $y\in L^\times$ satisfying
\begin{equation}\label{eq:pointwiseexpcondition}
 y^\dagger=\DD w
\end{equation}
if and only if $B\in\OO$, equivalently $w\in\mathcal S$.
In particular, no nonzero-valued global map $\mathcal E:L\to L^\times$
can satisfy
\[
 \DD\mathcal E(w)=\mathcal E(w)\DD w
 \quad\text{for every }w\in L.
\]
No homomorphism axiom is needed for this impossibility.
\end{theorem}
\begin{proof}
By \cref{thm:logderivativeimage}, condition
\eqref{eq:pointwiseexpcondition} is equivalent to $\DD B\in\II$.
But $\DD^{-1}(\II)=\OO$, as proved in
\cref{thm:I,cor:quotients}. Thus it is equivalent to $B$ being finite.
For such $w$, \cref{thm:stripexp} provides a solution. For infinite $B$,
no value at that single $w$ can meet the condition. Taking $w=i\omega$
already rules out a global map.
\end{proof}

This theorem is about the fixed Berarducci--Mantova derivation, acting on
values. It does not forbid a global surcomplex exponential defined by
other structural requirements. In particular, the ordered-exponential and
surcomplex exponential constructions studied by Ehrlich and Kaplan
\cite{EhrlichKaplan} should not be described as nonexistent. The
obstruction says that an extension to infinite imaginary arguments cannot
also satisfy \eqref{eq:stripcompatibility} for this $\DD$.

\begin{example}[Why discarding an infinite phase changes the calculus]
A convention that discards the purely infinite part of an angle can assign
phase $1$ to $\omega$. Such a value is algebraically meaningful under its
specified convention, but its intrinsic derivative is $\DD 1=0$.
Differential exponential compatibility at $i\omega$ would instead require
$\DD 1=i$. Thus a phase convention and a differential solution are
mathematically different objects, even when both are called
$e^{i\omega}$.
\end{example}

There is no contradiction with ordinary sine and cosine lifts on finite
halos. They obey the chain rule with the factor $\DD z$. For a finite
$z$, this factor is infinitesimal; there is no finite surreal $z$ with
$\DD z=1$. The forbidden oscillator would require a genuinely
infinite intrinsic clock.

\section{Nonoscillation and constant-coefficient equations}
\label{sec:constantODE}

\subsection{The elementary oscillator obstruction}
\begin{corollary}\label{cor:oscillator}
The equations
\[
 \DD y=iy,\qquad \DD y=-iy,\qquad \DD^2y+y=0
\]
have only the zero solution in $L$.
\end{corollary}
\begin{proof}
The first two equations would require $1\in\II$ or $-1\in\II$,
contradicting $\II\subseteq\mm$. For the third equation,
\[
 \DD^2+1=(\DD-i)(\DD+i).
\]
Both factors are injective by the first two cases, so their product is
injective.
\end{proof}

There is a complementary proof for a real solution $y$. Its energy
$y^2+(\DD y)^2$ has derivative zero, hence is an ordinary nonnegative
real constant. Therefore $y$ and $\DD y$ are finite. The derivative of
a finite surreal is infinitesimal, so $\DD y$ is infinitesimal, and
$\DD^2y=-y$ then forces $y$ infinitesimal too. The real energy constant
is infinitesimal and hence zero. Applying this to the real and imaginary
parts proves the complex case. The factorization proof, however, scales
more efficiently to general operators.

\subsection{All scalar constant-coefficient kernels}
For a nonzero polynomial $P(X)\in\C[X]$, the operator $P(\DD)$ is
$\C$-linear on $L$. Let $m_r$ be the multiplicity of a real root $r$ of
$P$.

\begin{theorem}[Exact constant-coefficient solution space]
\label{thm:constantoperators}
One has
\begin{equation}\label{eq:constantoperators}
 \ker P(\DD)=
 \bigoplus_{\substack{r\in\R\\P(r)=0}}
   e^{r\omega}\{q(\omega):q\in\C[X],\ \deg q<m_r\}.
\end{equation}
In particular, its complex dimension is the sum of the multiplicities of
\emph{real} roots of $P$, not necessarily $\deg P$.
\end{theorem}
\begin{proof}
First, if $\lambda\in\C\setminus\R$, its imaginary part is a
nonzero real constant and is not in $\II$. Thus $\DD-\lambda$ is
injective by \cref{thm:logderivativeimage}. Every power and every finite
product of such factors is injective.

Second, for real $r$ and every $u\in L$,
\[
 (\DD-r)(e^{r\omega}u)=e^{r\omega}\DD u.
\]
Iterating gives
\[
 (\DD-r)^m(e^{r\omega}u)=e^{r\omega}\DD^m u.
\]
We claim that $\ker\DD^m=\C[\omega]_{<m}$. This is clear for
$m=1$ since the constants are $\C$. If it holds for $m-1$ and
$\DD^m u=0$, then $\DD u$ is a polynomial in $\omega$ of degree
at most $m-2$. Integrate that polynomial termwise; the difference from
$u$ is a complex constant. This proves the claim by induction, and gives
the summand corresponding to $r$.

Finally, factor $P=cQR$, with $Q$ the product of its real-root factors
and $R$ the product of its nonreal-root factors. Since $R(\DD)$ is
injective and commutes with $Q(\DD)$,
$\ker P(\DD)=\ker Q(\DD)$. For pairwise coprime constant
polynomials, the usual Bezout identity decomposes the kernel of their
product as the direct sum of their kernels. Concretely, if $Q=AB$ and
$uA+vB=1$, then a vector $y\in\ker(AB)(\DD)$ decomposes as
$u(\DD)A(\DD)y+v(\DD)B(\DD)y$, in $\ker B(\DD)$ and
$\ker A(\DD)$ respectively. The intersection is zero by the same
identity. Repeating gives \eqref{eq:constantoperators}.
\end{proof}

For example,
\[
 \ker(\DD^2-1)=\C e^\omega\oplus\C e^{-\omega},\qquad
 \ker(\DD^2+1)=0,
\]
while
\[
 \ker\!\bigl(\DD(\DD-1)^2(\DD^2+1)\bigr)
     =\C\oplus e^\omega(\C+\C\omega).
\]
The latter fifth-order operator has a three-dimensional solution space.
There is no contradiction with the ordinary five-dimensional space of
complex analytic solutions as functions of a free variable. Its two
oscillatory modes are not elements of this differential field.

\subsection{Constant matrix systems}
\begin{theorem}\label{thm:matrixsystems}
Let $A\in M_n(\C)$, and consider
\begin{equation}\label{eq:matrixsystem}
 \DD Y=AY,\qquad Y\in L^n.
\end{equation}
Its solution space has complex dimension equal to the sum of the
algebraic multiplicities of the real eigenvalues of $A$. On a Jordan block
$rI+N$ of size $m$ with $r\in\R$, all solutions are
\begin{equation}\label{eq:jordansolutions}
 Y=e^{r\omega}
       \left(\sum_{k=0}^{m-1}\frac{\omega^kN^k}{k!}\right)c,
 \qquad c\in\C^m.
\end{equation}
On a Jordan block with nonreal eigenvalue, only the zero solution exists.
An invertible fundamental solution matrix over $L$ exists if and only if
all eigenvalues of $A$ are real.
\end{theorem}
\begin{proof}
An ordinary complex similarity transformation is constant for $\DD$, so
we may pass to Jordan form. On a nonreal block, the last coordinate solves
$(\DD-\lambda)y_m=0$ and is zero. Upward substitution then makes
all coordinates zero.

On a real block, the finite matrix polynomial
$B=\sum_{k=0}^{m-1}\omega^kN^k/k!$ satisfies $\DD B=NB$ and
has inverse $\sum_{k=0}^{m-1}(-\omega)^kN^k/k!$. Thus multiplication
by $e^{r\omega}B$ transforms a constant vector into a solution, and its
inverse transforms any solution into a constant vector. This proves
\eqref{eq:jordansolutions} and the dimension count. The block formula
provides an invertible fundamental matrix when all eigenvalues are real;
otherwise every solution is zero on some nonreal block, preventing
invertibility.
\end{proof}

For the real rotation matrix
$\left(\begin{smallmatrix}0&-1\\1&0\end{smallmatrix}\right)$ the
solution space is zero. By contrast, a unipotent Jordan block has a
fundamental matrix that is a polynomial in $\omega$. This contrast
isolates oscillation rather than general complexity as the obstruction.

\subsection{What these kernel results do not say}
An operator can have zero kernel and still be surjective. Conversely,
knowing a full kernel does not settle an inhomogeneous equation over a
specified smaller subfield. The preceding theorems classify homogeneous
solutions in $L$ only. The next section gives an explicit field on which
$\DD-i$ is bijective. No claim of differential closedness, universal
Picard--Vessiot closure, or a decision procedure for arbitrary nonlinear
differential equations is being made.

\section{An exact computational model: Laurent series at infinity}
\label{sec:laurent}

\subsection{A derivation-stable but not integration-closed field}
Put $t=\omega^{-1}$ and identify
\[
 F=\C((t))\subset L
\]
by Hahn evaluation. Its elements are formal Laurent series
$f=\sum_{n\ge n_0}c_nt^n$, with $n_0\in\Z$ and $c_n\in\C$.
The intrinsic derivation is exactly
\begin{equation}\label{eq:LaurentD}
 D=-t^2\frac{d}{dt},\qquad
 Df=\sum_{n\ge n_0}(-n)c_nt^{n+1}.
\end{equation}
Indeed, $\DD t=-t^2$, and strong additivity gives the formula for all
Laurent series. We write $D$ here to emphasize the concrete operator,
although it is the restriction of $\DD$.

\begin{proposition}\label{prop:Laurentintegration}
The constants of $(F,D)$ are $\C$, and
\begin{equation}\label{eq:Laurentimage}
 D(F)=\{f\in F:[t^1]f=0\}.
\end{equation}
For $f=\sum c_nt^n$ with $c_1=0$, a primitive is
\begin{equation}\label{eq:Laurentprimitive}
 \sum_{n\ne1}\frac{c_n}{1-n}t^{n-1}.
\end{equation}
\end{proposition}
\begin{proof}
A coefficient of $t^1$ in $Df$ could only come from differentiating
$t^0$, whose derivative vanishes. All other powers have the primitives
shown in \eqref{eq:Laurentprimitive}; shifting the lower exponent by one
preserves Laurent support. The same coefficient calculation shows that
$Df=0$ forces all coefficients except $c_0$ to vanish.
\end{proof}

Thus $t$ has no primitive in $F$, although it has the primitive
$\log\omega$ in $K$. This is not a failure of ambient surjectivity; it
is a failure of closure of a particular representation parent.

For another comparison, the power-series ideal $t^2\C[[t]]$ is the
image under $D$ of $t\C[[t]]$. On the real Laurent field this is precisely
$\II\cap\R((t))$: a finite Laurent primitive requires the absence of
all powers below $t^2$, and conversely the displayed primitive is
infinitesimal. A primitive outside the Laurent field cannot rescue a
missing term, because primitives differ only by real constants. More
explicitly, subtract the Laurent primitive of all terms except $c_1t$;
the remaining primitive is $c_1\log\omega$. Any nonzero $c_1$, or any
nonzero term whose primitive has a negative power of $t$, prevents
finiteness. The logarithmic term cannot cancel a Laurent principal part.

\subsection{Resolvents with a support certificate}
\begin{theorem}[Exact Laurent resolvent]\label{thm:resolvent}
For every $\lambda\in\C^\times$, the operator $D-\lambda$ is a
bijection of $F$. Its inverse is the strongly summable operator series
\begin{equation}\label{eq:resolvent}
 R_\lambda f=-\sum_{j\ge0}\lambda^{-(j+1)}D^j f.
\end{equation}
For its finite truncation
$R_{\lambda,N}f=-\sum_{j=0}^N\lambda^{-(j+1)}D^j f$, the residual is
exactly
\begin{equation}\label{eq:resolventresidual}
 (D-\lambda)R_{\lambda,N}f
   =f-\lambda^{-(N+1)}D^{N+1}f.
\end{equation}
\end{theorem}
\begin{proof}
Write $f=p+g$, where $p\in\C[t^{-1}]$ is the finite part supported at
nonpositive powers and $g\in t\C[[t]]$. Repeated application of $D$
to $p$ eventually gives zero: under $t^{-1}=\omega$ this is ordinary
differentiation of a polynomial in $\omega$. For $g$, every nonzero
application of $D$ raises its least exponent by at least one. Hence the
family $(D^jg)_j$ has a well-ordered union of exponents and only finitely
many contributions to any fixed exponent. The finite family coming from
$p$ causes no problem. This proves strong summability of
\eqref{eq:resolvent}.

The finite identity \eqref{eq:resolventresidual} follows by cancellation
of adjacent terms. Direct coefficientwise cancellation in the summable
infinite family then gives $(D-\lambda)R_\lambda f=f$. Equivalently,
one can use strong additivity; there is no fine-topological limit.

For injectivity, let $y\ne0$ have least exponent $n$. If $n\ne0$, the
least exponent of $Dy$ is $n+1$, so the coefficient of $t^n$ in
$(D-\lambda)y$ is $-\lambda[t^n]y\ne0$. If $n=0$, the constant
term of $Dy$ is zero and the constant term of $(D-\lambda)y$ is again
nonzero. Thus its kernel in $F$ is zero.
\end{proof}

For real nonzero $\lambda$, the ambient equation
$(\DD-\lambda)y=0$ has nonzero solutions $ce^{\lambda\omega}$,
but none lies in $F$. For nonreal $\lambda$, it has none even in $L$.
The same Laurent inverse exists in both cases. Parent selection affects
homogeneous solution spaces and must remain part of an algorithm's
contract.

\subsection{A factorially divergent but exact solution}
For $f=t$ and $\lambda=i$,
\[
 D^jt=(-1)^j j!t^{j+1}.
\]
Substitution in \eqref{eq:resolvent} gives the exact Hahn series
\begin{equation}\label{eq:factorialsolution}
 y=\sum_{j\ge0}i^{j+1}j!t^{j+1}
   =it-t^2-2it^3+6t^4+24it^5-120t^6+\cdots,
 \qquad Dy-iy=t.
\end{equation}
Its corresponding ordinary power series has radius of convergence zero,
since its coefficient magnitudes are factorials. Yet its support is the
positive integer lattice, so it is a perfectly valid surcomplex value.
This example demonstrates why analytic convergence, formal summability,
and exact differential solvability must not be used interchangeably.

Let $y_N$ be the sum through $j=N$. The exact residual is
\begin{equation}\label{eq:factorialresidual}
 (D-i)y_N-t
   =-i^{-(N+1)}(-1)^{N+1}(N+1)!t^{N+2}.
\end{equation}
This is a useful certificate: a finite program can verify it without
representing the entire infinite series. It certifies a formal remainder
at a stated order, not convergence of the sequence $(y_N)$ in the full
surreal topology.

\subsection{A coefficient algorithm and its precise guarantee}
Writing $(D-\lambda)y=f$ coefficientwise yields
\begin{equation}\label{eq:coefficientrecursion}
 -\lambda y_n-(n-1)y_{n-1}=f_n,
 \qquad
 y_n=-\lambda^{-1}\bigl(f_n+(n-1)y_{n-1}\bigr).
\end{equation}
Below a known lower exponent bound, all coefficients are zero. The
recursion computes each requested finite prefix with one field inversion
of $\lambda$ and a linear number of coefficient operations. It works
also for negative lower bounds. It does not make equality of arbitrary
program-generated streams decidable.

For a nonzero polynomial $P\in\C[X]$ with $P(0)\ne0$, factor it
over $\C$ and compose the resolvents, with multiplicities. This proves
that $P(D)$ is a bijection on $F$. The proof is algebraic; an effective
implementation additionally needs exact representations of the constant
roots, or can derive a direct coefficient recurrence from $P$ instead.
When $P(0)=0$, the factor $D$ brings the obstruction in
\eqref{eq:Laurentimage}; it must not be handled by division by zero in a
resolvent routine.

\section{Set-sized differential workspaces}
\label{sec:workspaces}

The repository's field-localization principle explains why a set of
surcomplex data can be handled inside a set-sized algebraic workspace.
Differential tasks need additional closure operations. The following
existence result states exactly what can be arranged without claiming an
effective presentation or a closure under all Hahn sums.

\begin{theorem}[Differential and local-analytic workspace]
\label{thm:workspace}
Let $S\subseteq K$ be a set. There exists a set-sized real-closed
subfield $K_0\subseteq K$ containing $S\cup\R\cup\{\omega\}$
and closed under $\DD$, $\Prim$, surreal exponential, and logarithm
of positive elements. It can also be chosen so that, for every
$h\in K_0(i)\cap\mm_L$ and every $f\in\C[[X]]$, the Hahn value
$f(h)$ belongs to $K_0(i)$.

The cardinality can be bounded by
$\kappa=\max(|S|,|\R|,\aleph_0)$. The complexification $K_0(i)$ is
algebraically closed, has differential constants exactly $\C$, and is
closed under a chosen primitive of every element. For $a,b\in K_0$, the
ambient criterion $b\in\II$ for a nonzero solution of
$\DD y=(a+ib)y$ is already realized inside $K_0(i)$.
\end{theorem}
\begin{proof}
Start with the field generated by $S\cup\R\cup\{\omega\}$.
Given a set-sized intermediate field $E_n$, adjoin all values
\[
 \DD x,\quad \Prim(x),\quad \exp x\quad(x\in E_n),
 \qquad \log x\quad(x\in E_n,\ x>0).
\]
Also adjoin the real and imaginary parts of all values $f(h)$ with
$f\in\C[[X]]$ and $h\in E_n(i)\cap\mm_L$. Finally take the
relative real algebraic closure of the generated field in $K$ to obtain
$E_{n+1}$. Class replacement ensures that the image of each of these
specified sets under the indicated class operation is a set.

There are $|\R|$ formal complex power series, since
$|\C|^{\aleph_0}=|\R|$. At each stage at most $\kappa$ elements are
adjoined, and an algebraic closure of a field of that cardinality has the
same cardinality. Thus $K_0=\bigcup_{n<\omega}E_n$ is a set of
cardinality at most $\kappa$. Every operation applied to an element of
$K_0$ is included at the next stage containing that element. Every finite
polynomial with coefficients in $K_0$ occurs at one stage, so $K_0$ is real
closed. The formal-series closure follows in the same manner.

Since $K_0\subset K$ contains all reals, its constant field is exactly
$\R$, and $K_0(i)$ has constants $\C$. Closure under $\Prim$ yields
surjective differentiation on $K_0$, and then on $K_0(i)$. Real closedness
implies algebraic closedness of $K_0(i)$.

Finally, if $b\in K_0\cap\II$, then $\Prim(b)=J(b)\in K_0\cap\mm$.
The explicit solution
$\exp(\Prim(a))\Exp(iJ(b))$ therefore belongs to $K_0(i)$ by its
closure properties. If $b\notin\II$, there is no nonzero solution
in the ambient field, hence none in $K_0(i)$.
\end{proof}

In fact, every nonzero solution of the homogeneous first-order equation
with coefficients in $K_0(i)$ belongs to $K_0(i)$ whenever one exists:
its multiplier is in $\C$, which was included from the beginning.
The same holds for the constant-coefficient solution spaces in
\cref{thm:constantoperators,thm:matrixsystems}.

\begin{remark}[Closure does not mean Hahn completeness]
The theorem does not assert that $K_0$ contains every Hahn sum of every
summable family of its elements. The added formal evaluations have
constant coefficients and a \emph{single infinitesimal argument}; they
are a specified set of operations. Unrestricted summation closure is a
much stronger requirement. A particular coherent family and its requested
evaluations can be included as set-sized data, but arbitrary future
families are not automatically covered by the theorem.
\end{remark}

The workspace proof is an existence construction by repeated closure. It
has no birthday bound, no computable enumeration, and no termination
claim for symbolic integration. More specialized results about surreal
fields stable under exponential, logarithmic, derivative and
antiderivative operations are available in the literature
\cite{BournezGuilmant}; the elementary closure argument above is sufficient
for the finite-data foundational issue addressed here.

\section{Composition boundaries and implementation contracts}
\label{sec:contracts}

\subsection{Evaluation is not unrestricted composition}
The total-derivative theorem differentiates a specified common-domain
Hahn family at a finite argument. It should not be restated as a
composition operation on all surreal numbers.

Berarducci and Mantova construct a composition theory for their
$\omega$-series and prove a chain rule there. They also prove that the
Berarducci--Mantova derivation on \emph{all} of $\No$ admits no
composition compatible with their stated axioms
\cite[Theorem 8.4]{BMgerms}. The hypotheses matter: their compatibility
includes the appropriate constant rule, order behaviour for positive
derivative, and a chain rule at positive infinite arguments. The theorem
is not a claim that every conceivable notion of composition, or every
different choice of derivation, is impossible.

Thus three operations must remain distinct: formal composition of
constant-coefficient series at infinitesimals, analytic evaluation of a
common-domain family, and composition of transseries-like elements at
positive infinite arguments. The first two were proved here with their
support certificates. The third requires a separate function class and
its own theorems.

\subsection{A minimal derivation-aware representation interface}
A useful exact parent should record both its ambient embedding and its
closure capabilities. For example, a Laurent parent may advertise
\[
 \text{closed under }D,\quad
 \text{closed under }(D-\lambda)^{-1}\ (\lambda\ne0),\quad
 \text{not closed under primitives}.
\]
A log-exponential extension can support more primitives without thereby
supporting arbitrary global phases. Adding a constructor named
\texttt{Exp} is not enough: its domain and derivative contract must be
stated.

\begin{center}
\begin{tabular}{@{}>{\raggedright\arraybackslash}p{.23\textwidth}>{\raggedright\arraybackslash}p{.37\textwidth}>{\raggedright\arraybackslash}p{.31\textwidth}@{}}
\toprule
Operation & Required certificate & Failure or limitation \\
\midrule
Intrinsic derivative & Chosen derivation; supported image in a known parent & A formal monomial derivative need not lie in the old parent \\
Hahn evaluation & Joint reverse-well-ordered support and finite multiplicity & A formal expression may not denote a sum at the input \\
Coherent total derivative & One common coefficient domain; both derivatives & Dropping $\Dc$ changes the answer \\
Exponential integration & A finite primitive of the imaginary coefficient & Smallness alone is insufficient \\
Laurent resolvent & $\lambda\ne0$ and a lower exponent bound & Prefix computation is not stream equality \\
Global phase & Explicit phase convention and supported domain & A group law does not imply BM compatibility \\
\bottomrule
\end{tabular}
\end{center}

For a practical implementation, the first reliable targets are rational
functions in $\omega$, Laurent streams with exact coefficient access,
and finite towers of explicitly represented logarithms and exponentials.
At each extension, a derivative routine should return both a value and
the parent in which it is represented. An integration routine should
separate a proved primitive, a proved obstruction in the selected parent,
and a failure of the available algorithm to decide. The Laurent missing
$t$-coefficient is a proved parent obstruction; failure to recognize a
complicated primitive in the full surreal field is not.

\subsection{What a verification record can establish}
The accompanying program treats $t$ and $\omega$ as formal variables in
exact symbolic algebra. It checks finite product and chain identities,
Laurent coefficient recurrences, resolvent residuals, truncated small
exponential/logarithm identities, and examples of scalar and matrix
differential equations. These checks detect sign, index and truncation
errors. They do not verify the BM construction, class foundations,
Hahn--Neumann lemma, or infinite-support arguments. Those remain the
mathematical proofs and external inputs explicitly identified above.

Numerical substitution $t=10^{-6}$ would not verify a Hahn identity and
is not used. Nor does the program treat a successful truncation as a
convergence theorem. Its output is an exact finite certificate at the
specified order.

\section{Synthesis and placement in the documentation}
\label{sec:synthesis}

The missing subject is not another list of formal derivative rules. It is
the interaction among a derivation on numbers, a derivative in a function
variable, and the restricted phase geometry of a nonoscillatory ordered
differential field.

The analytic interface is the commuting pair $(\Dc,\Dz)$ and its
supported evaluation rule. The differential-algebraic interface is the
ideal $\II=\DD(\mm)$, which measures exactly which imaginary
logarithmic derivatives can occur. Together they explain why a rich
surcomplex analytic calculus can coexist with the absence of a nonzero
intrinsic oscillator. They also explain why a global algebraic phase
extension and a derivative-compatible exponential are different promises.

A suitable home for this article is a new package:
\begin{quote}\pathtext{docs/surreal/intrinsic-differential-calculus/}\end{quote}
It should have cross-links
from the analysis, trigonometry, and computer-algebra reports. The current
phase constructions need not be removed or renamed as false; their
contracts need only be distinguished from the intrinsic differential one.
The derivation itself should be named wherever it is used.

The most useful next formalization target is modest and sharply bounded:
start with a real-closed exponential differential field carrying the
stated strong summation operations, formalize the infinitesimal polar
lemma, and derive the logarithmic-derivative image theorem. The difficult
BM existence construction can remain an explicit imported interface
until separately formalized. This is a suggested decomposition of proof
obligations, not a claim that Lean code is included or that the article
has been machine checked.

\appendix
\section{Notation and hypotheses at a glance}
\label{app:notation}
\begin{longtable}{@{}>{\raggedright\arraybackslash}p{.19\textwidth}>{\raggedright\arraybackslash}p{.72\textwidth}@{}}
\toprule
Symbol & Meaning and domain \\
\midrule
\endfirsthead
\toprule
Symbol & Meaning and domain \\
\midrule
\endhead
$K,L$ & $\No$ and $\No[i]$, treated as class fields; individual supports and computational workspaces are sets. \\
$\DD$ & The fixed Berarducci--Mantova derivation, normalized by $\DD\omega=1$, extended to $L$. \\
$\Dz$ & Analytic differentiation in the variable $Z$; coefficients in $L$ are held fixed. \\
$\Dc$ & Intrinsic differentiation of the Hahn coefficients/monomials of a common-domain family; the ordinary variable is held fixed. \\
$\OO,\mm$ & The finite surreals and the infinitesimals, with $\OO=\R\oplus\mm$. \\
$\st$ & Standard part on finite elements only; a ring homomorphism. \\
$\ct$ & The coefficient of the Conway monomial $1$, defined on all $K$; real-linear but not multiplicative. \\
$\Prim$ & The primitive normalized by $\ct\Prim(b)=0$; defined on all $K$ using surjectivity. \\
$\II$ & $\DD(\OO)=\DD(\mm)$, a nonzero convex ideal of $\OO$ strictly smaller than $\mm$. \\
$J$ & The inverse $\II\to\mm$ of the restriction of $\DD$; the unique infinitesimal primitive. \\
$z^\dagger$ & $(\DD z)/z$ for $z\ne0$; it does not require a preselected logarithm. \\
$\Exp,\Log$ & Mutually inverse formal functions on $\mm_L$ and $1+\mm_L$. \\
$E$ & The derivative-compatible exponential on $\mathcal S=K+i\OO$; $K_0$ denotes a set-sized workspace field. \\
$t,D$ & $t=\omega^{-1}$ and $D=-t^2d/dt$ on $\C((t))$. \\
$\omega^a$ & Conway monomial notation; for general surreal $a$ it is not a generic exp--log power constructor. \\
\bottomrule
\end{longtable}

\section{Logical dependencies and reproducibility}
\label{app:dependencies}

\subsection{Dependency map}
The ideal theorem uses only the constant field, surjective derivation,
$H$-field sign axiom, and the normalization $\DD\omega=1$ for the
smallness proof. The polar decomposition additionally uses real
closedness, standard part, and infinitesimal formal exponential/logarithm.
The logarithmic-derivative theorem then uses real exponential
compatibility. Constant-coefficient kernel classification is finite
linear differential algebra once that theorem is established.

Strong additivity is essential for the infinitesimal chain rule and the
coherence-preservation theorem. It is not supplied by the Leibniz rule
alone. Conversely, the Laurent identities can be verified within formal
series algebra independently of a construction of all surreals; the
embedding into $L$ is what identifies their operator with $\DD$.

The article does not prove the uniqueness of the chosen surreal
derivation, unrestricted closure under composition, differential
closedness of $L$, or the existence of a universally effective CAS for
surreal numbers. Its maximal-strip theorem does not deny the existence
of global phase maps with different axioms.

\subsection{Build and check}
The source is standalone, with an embedded bibliography and no external
figures. From the package directory:
\begin{verbatim}
latexmk -pdf -interaction=nonstopmode -halt-on-error article.tex
python code/verify_identities.py --output data/verification.json
\end{verbatim}
The delivered run used Python 3.13.5 and SymPy 1.14.0 and passed
259 of 259 checks. A matrix identity or a whole checked coefficient
prefix counts as one check. These numbers describe finite symbolic
identities, not 259 machine-checked theorems about the full surreal field.
The report records every result. The accompanying README records the
final PDF page count and build status. TeX auxiliary files are not part
of the delivered archive.

\begin{thebibliography}{99}
\raggedright

\bibitem{RepoMap}
V.~Reshetnikov, \emph{The surreal and surcomplex reports}, repository
navigation document, commit
\texttt{e260237db9b71da8b74a0c13c8e6355119091100}.
\url{https://github.com/VladimirReshetnikov/Surreal/blob/e260237db9b71da8b74a0c13c8e6355119091100/docs/README.md}.

\bibitem{RepoManifest}
V.~Reshetnikov, \emph{The surreal and surcomplex reports}, typeset
catalogue source, same repository commit.
\url{https://github.com/VladimirReshetnikov/Surreal/blob/e260237db9b71da8b74a0c13c8e6355119091100/docs/manifest.tex}.

\bibitem{RepoAnalysis}
\emph{Surcomplex analysis: holomorphic functions on $K=\No[i]$}, package
scope and reading guide, same repository commit. AI-assisted research
drafts; cited for their stated scope, not as independently verified theorems.
\url{https://github.com/VladimirReshetnikov/Surreal/blob/e260237db9b71da8b74a0c13c8e6355119091100/docs/surcomplex/analysis/README.md}.

\bibitem{RepoTrig}
\emph{Trigonometry on the Surcomplex Plane: Canonical Phases,
Arbitrary-Scale Geometry, and Infinitesimal Degeneration}, package scope
and explicit non-claims, same repository commit.
\url{https://github.com/VladimirReshetnikov/Surreal/blob/e260237db9b71da8b74a0c13c8e6355119091100/docs/surcomplex/trigonometry/README.md}.

\bibitem{RepoCAS}
\emph{Surreal and Surcomplex Numbers in Symbolic Computer Algebra},
package scope and capability discussion, same repository commit.
\url{https://github.com/VladimirReshetnikov/Surreal/blob/e260237db9b71da8b74a0c13c8e6355119091100/docs/foundations-and-computation/computer-algebra/README.md}.

\bibitem{BM}
A.~Berarducci and V.~Mantova,
\emph{Surreal numbers, derivations and transseries},
arXiv:1503.00315v3, 2015. Theorem numbering cited in the text refers to
this version.
\url{https://arxiv.org/abs/1503.00315v3}.

\bibitem{BMpub}
A.~Berarducci and V.~Mantova,
\emph{Surreal numbers, derivations and transseries},
Journal of the European Mathematical Society \textbf{20} (2018),
339--390.
\url{https://doi.org/10.4171/JEMS/769}.
This is the published version of \cite{BM}, not an independent source.

\bibitem{ADHsurreal}
M.~Aschenbrenner, L.~van den Dries and J.~van der Hoeven,
\emph{The surreal numbers as a universal $H$-field},
arXiv:1512.02267v3, 2016.
\url{https://arxiv.org/abs/1512.02267v3}.

\bibitem{BMgerms}
A.~Berarducci and V.~Mantova,
\emph{Transseries as germs of surreal functions},
Transactions of the American Mathematical Society \textbf{371} (2019),
3549--3592; arXiv:1703.01995v3, 2017.
\url{https://doi.org/10.1090/tran/7428};
\url{https://arxiv.org/abs/1703.01995v3}.
The composition obstruction is Theorem~8.4 in the cited arXiv version.

\bibitem{EhrlichKaplan}
P.~Ehrlich and E.~Kaplan,
\emph{Surreal ordered exponential fields},
arXiv:2002.07739v3, 2021.
\url{https://arxiv.org/abs/2002.07739v3}.

\bibitem{BournezGuilmant}
O.~Bournez and Q.~Guilmant,
\emph{Surreal fields stable under exponential, logarithmic, derivative
and anti-derivative functions}, arXiv:2211.08396, 2022.
\url{https://arxiv.org/abs/2211.08396}.

\end{thebibliography}
\end{document}
